\chapter{Nutritional and Metabolic Disorders}
\label{ch:nutrition}
Nutritional and metabolic disorders encompass a wide spectrum of diseases resulting from deficiencies or excesses of dietary nutrients, impaired metabolism due to enzyme defects, and dysregulation of energy homeostasis. These conditions range from global public health problems affecting billions---such as protein-energy malnutrition, iron deficiency anemia, and iodine deficiency---to rare inborn errors of metabolism with incidences of 1 in 100,000 or fewer. Their manifestations affect virtually every organ system, and early recognition is critical to prevent irreversible sequelae. Nutritional biochemistry and metabolic pathways form the foundation for understanding these disorders, linking dietary intake to cellular function, growth, and development.

%-------------------------------------------------------------
\section{Disorders of Lipid Metabolism}
%-------------------------------------------------------------

%-------------------------------------------------------------

%-------------------------------------------------------------

%-------------------------------------------------------------

%-------------------------------------------------------------

%-------------------------------------------------------------

\label{sec:Disorders of Lipid Metabolism}
%-------------------------------------------------------------%-------------------------------------------------------------

\subsection{Dysbetalipoproteinemia}
\label{sec:Dysbetalipoproteinemia}
Dysbetalipoproteinemia (type III hyperlipoproteinemia, remnant removal disease) is a disorder of chylomicron and VLDL remnant clearance caused by homozygosity for apolipoprotein E2 (apoE2/E2 genotype) in conjunction with a second metabolic hit (obesity, diabetes, hypothyroidism, or estrogen deficiency). Only 1--5\% of individuals with the apoE2/E2 genotype develop the full phenotype, which is otherwise rare (prevalence approximately 1 in 10,000). The condition results from impaired hepatic clearance of remnant lipoproteins due to poor binding of apoE2 to the LDL receptor. Clinical features include both cholesterol and triglyceride elevation (typically in equal proportion, greater than 300 mg/dL), tuberoeruptive xanthomas, and palmar crease xanthomas (pathognomonic). Diagnosis is by ultracentrifugation showing a broad $\beta$-band, VLDL cholesterol-to-total triglyceride ratio greater than 0.30, and apoE genotyping (E2/E2). Management includes lifestyle modification; statins and fibrates are synergistic and often normalize lipids. Prognosis is favorable with treatment.

\subsection{Familial Combined Hyperlipidemia}
\label{sec:Familial Combined Hyperlipidemia}
Familial combined hyperlipidemia (FCHL) is the most common genetic dyslipidemia, with a prevalence of 1--2\% in the general population and up to 10--20\% in patients with premature coronary artery disease. It is a complex oligogenic disorder characterized by overproduction of apolipoprotein B-100 and very-low-density lipoprotein (VLDL) by the liver. Clinical features include premature coronary artery disease, obesity, insulin resistance, and the absence of tendon xanthomas (distinguishing it from FH). The lipid phenotype is variable: affected individuals may have elevated total cholesterol and/or elevated triglycerides. Diagnosis is made in individuals with elevated apolipoprotein B (greater than 120 mg/dL), elevated LDL-C and/or triglycerides, and a family history of premature cardiovascular disease. Management includes lifestyle modification, statins as first-line therapy, with the addition of fibrates for hypertriglyceridemia. Prognosis is favorable with treatment.

\subsection{Familial Hypercholesterolemia}
\label{sec:Familial Hypercholesterolemia}
Familial hypercholesterolemia (FH) is an autosomal codominant disorder of LDL metabolism caused by loss-of-function mutations in the LDLR gene (LDL receptor, 90\% of cases), APOB (apolipoprotein B-100), or gain-of-function mutations in PCSK9. It is the most common monogenic lipid disorder, with a prevalence of 1 in 250 for heterozygous FH (HeFH) and 1 in 300,000 for homozygous FH (HoFH). The condition results from impaired hepatic clearance of LDL particles. Clinical features include tendon xanthomas (Achilles, extensor tendons of the hands), tuberous xanthomas, xanthelasmas, corneal arcus before age 45, and premature atherosclerotic cardiovascular disease; homozygous FH presents in childhood with LDL-C greater than 500 mg/dL, widespread xanthomas, and coronary atherosclerosis before age 20. Diagnosis is by clinical criteria (Dutch Lipid Clinic Network or Simon Broome criteria) supported by genetic testing. Management includes high-dose statins with ezetimibe as first-line for HeFH; PCSK9 inhibitors for inadequate response; for HoFH, all available agents and LDL apheresis are needed. Prognosis is excellent with early diagnosis and treatment.

\subsection{Lipoprotein(a) Disorders}
\label{sec:Lipoprotein(a) Disorders}
Lipoprotein(a) (Lp(a)) disorder is a lipid disorder defined by elevated Lp(a) levels (greater than 50 mg/dL or greater than 125 nmol/L), which is an independent, causal risk factor for atherosclerotic cardiovascular disease and calcific aortic valve stenosis. Lp(a) levels are predominantly genetically determined (by the LPA gene) and are largely unaffected by diet or lifestyle; the prevalence of elevated Lp(a) is approximately 20\% of the population worldwide. The condition results from apolipoprotein(a) component sharing homology with plasminogen and having antifibrinolytic properties, also promoting vascular inflammation and foam cell formation. Clinical features include premature cardiovascular disease and a family history of early heart disease or aortic stenosis, with no other specific clinical findings. Diagnosis is by measurement of Lp(a) mass or molar concentration (once in a lifetime, as levels are stable). Standard lipid-lowering therapies do not significantly lower Lp(a); niacin and PCSK9 inhibitors produce modest reductions; investigational therapies (pelacarsen, olpasiran) reduce Lp(a) by 80--90\%. Management focuses on aggressive control of other cardiovascular risk factors.

%-------------------------------------------------------------
\section{Inborn Errors of Metabolism}
%-------------------------------------------------------------

%-------------------------------------------------------------

%-------------------------------------------------------------

%-------------------------------------------------------------

%-------------------------------------------------------------

%-------------------------------------------------------------

\label{sec:Inborn Errors of Metabolism}
%-------------------------------------------------------------%-------------------------------------------------------------

\subsection{Alkaptonuria --- Ochronosis}
\label{sec:Alkaptonuria}
Alkaptonuria is a rare autosomal recessive disorder caused by deficiency of homogentisate 1,2-dioxygenase (HGD), the enzyme responsible for the breakdown of homogentisic acid (HGA) in the tyrosine degradation pathway. The incidence is approximately 1 in 250,000--1,000,000. The condition results from accumulation of HGA which polymerizes to form a dark brown-black pigment that deposits in connective tissues throughout the body. Clinical features include homogentisic aciduria (urine turns dark brown/black upon exposure to air), ochronosis (bluish-black pigmentation of ears, nose, sclera, and skin), ochronotic arthropathy (progressive degenerative joint disease affecting the spine and large joints), and cardiovascular involvement including aortic valve stenosis and coronary artery calcification. Diagnosis is by elevated urinary homogentisic acid and HGD gene sequencing. Management is symptomatic: nitisinone reduces HGA production and slows disease progression but does not reverse existing ochronosis; pain management, physical therapy, and joint replacement surgery for advanced arthropathy. Prognosis is variable; the disease is progressive.

\subsection{Alpha-1 Antitrypsin Deficiency}
\label{sec:Alpha-1 Antitrypsin Deficiency}
Alpha-1 antitrypsin (AAT) deficiency is an autosomal codominant disorder caused by mutations in the SERPINA1 gene, most commonly the Pi*Z allele, leading to reduced serum levels and accumulation of misfolded AAT in the liver. The incidence is approximately 1 in 2,500--5,000 live births in Northern European populations. The condition results from accumulation of polymerized AAT in hepatocytes causing liver disease and lack of functional AAT in the lungs permitting unchecked proteolytic destruction of alveolar walls. Clinical features include neonatal cholestasis, hepatitis, cirrhosis, and hepatocellular carcinoma (liver disease); early-onset panacinar emphysema particularly in smokers (lung disease); and widespread panniculitis. Diagnosis is by low serum AAT level, phenotyping (Pi*ZZ, Pi*SZ), and genotyping. Management includes intravenous augmentation therapy (weekly infusions of pooled human plasma-derived AAT) to slow lung function decline; smoking cessation is critical; liver transplantation is curative for end-stage liver disease. Prognosis is variable; avoidance of hepatotoxins is essential.

Disorders of lipid metabolism encompass both primary (genetic) and secondary dyslipidemias that contribute substantially to the global burden of atherosclerotic cardiovascular disease. Genetic dyslipidemias result from mutations affecting lipoprotein production, processing, or clearance, while secondary causes include diabetes, hypothyroidism, nephrotic syndrome, cholestasis, and medications. Early identification and aggressive lipid-lowering therapy can substantially modify cardiovascular risk.

\subsection{Galactosemia}
\label{sec:Galactosemia}
Galactosemia is an autosomal recessive disorder of galactose metabolism caused most commonly by deficiency of galactose-1-phosphate uridylyltransferase (GALT), with an incidence of approximately 1 in 30,000--60,000 live births. The condition results from accumulation of galactose-1-phosphate and galactitol when galactose from dietary lactose cannot be converted to glucose. Clinical features present after the initiation of milk feeding with poor feeding, vomiting, diarrhea, failure to thrive, hypoglycemia, jaundice, hepatomegaly, ascites, and liver dysfunction; Escherichia coli neonatal sepsis is a frequent and life-threatening complication; cataracts develop due to galactitol accumulation in the lens. Diagnosis is by newborn screening (elevated total galactose), red blood cell GALT enzyme activity (absent or severely reduced), and genetic testing of the GALT gene. Management involves immediate elimination of dietary lactose and galactose (soy-based or elemental formula) for life. Prognosis is guarded despite treatment; long-term complications including speech apraxia, cognitive deficits, ataxia, tremor, and premature ovarian insufficiency are common.

\subsection{Glycogen Storage Diseases}
\label{sec:Glycogen Storage Diseases}
Glycogen storage diseases (GSDs) are a group of autosomal recessive disorders (except GSD IX, X-linked) resulting from defects in glycogen synthesis, degradation, or glycolysis, leading to abnormal glycogen accumulation. Von Gierke disease (GSD I) is caused by glucose-6-phosphatase deficiency, presenting with severe fasting hypoglycemia, hepatomegaly, growth retardation, lactic acidosis, hyperuricemia, hypertriglyceridemia, and a rounded, doll-like facies; treatment involves continuous nasogastric infusion of uncooked cornstarch. Pompe disease (GSD II) is caused by acid maltase deficiency, a lysosomal storage disorder presenting in infants with hypertrophic cardiomyopathy, macroglossia, hypotonia (floppy baby), and rapidly progressive weakness; enzyme replacement therapy (alglucosidase alfa) improves survival. Cori disease (GSD III) is caused by debranching enzyme deficiency, causing hypoglycemia, hepatomegaly, myopathy, and possible cardiomyopathy. McArdle disease (GSD V) is myophosphorylase deficiency (see Chapter~16, Muscles).

\subsection{Hereditary Hemochromatosis}
\label{sec:nutrition:Hereditary Hemochromatosis}
Hereditary hemochromatosis (HH) is an autosomal recessive disorder of iron metabolism characterized by excessive intestinal absorption of dietary iron, leading to progressive iron overload and organ damage. The most common form (HFE-related, type 1) is caused by the C282Y mutation in the HFE gene, with a prevalence of approximately 1 in 200--300 in Northern European populations. The condition results from the mutation impairing hepcidin production, leading to unregulated ferroportin-mediated iron efflux from enterocytes and macrophages. Clinical features typically manifest in adulthood and include hepatic cirrhosis (most common cause of morbidity), hepatocellular carcinoma (30\% risk with cirrhosis), skin hyperpigmentation (bronze diabetes), diabetes mellitus, hypogonadism, arthropathy (chondrocalcinosis), and cardiomyopathy. Diagnosis is by elevated transferrin saturation (greater than 45\%) and ferritin, confirmatory genetic testing for HFE mutations, and liver biopsy when indicated. Management is therapeutic phlebotomy (weekly until ferritin less than 50 ng/mL, then maintenance); iron chelation is used if phlebotomy is contraindicated. Prognosis is excellent with early diagnosis and treatment.

\subsection{Homocystinuria}
\label{sec:Homocystinuria}
Homocystinuria (classic form) is an autosomal recessive disorder caused by deficiency of cystathionine $\beta$-synthase (CBS), the enzyme that converts homocysteine to cystathionine in the transsulfuration pathway. The incidence is approximately 1 in 200,000--300,000 live births. The condition results from accumulation of homocysteine and methionine, producing a characteristic clinical phenotype resembling Marfan syndrome. Clinical features include tall, thin build, ectopia lentis (downward lens dislocation), myopia, intellectual disability (variable), skeletal abnormalities (scoliosis, pectus excavatum, arachnodactyly, genu valgum), osteoporosis, and a high risk of thromboembolism (the most life-threatening feature). Diagnosis is by elevated total plasma homocysteine (greater than 100 \textmu mol/L), elevated methionine, and low cystine; newborn screening enables early detection. Management involves pyridoxine for responsive patients; pyridoxine-nonresponsive patients require methionine restriction, betaine supplementation, and B12/folate supplementation; antiplatelet therapy reduces thromboembolic risk. Prognosis is favorable with early treatment.

\subsection{Lysosomal Storage Diseases}
\label{sec:Lysosomal Storage Diseases}
Lysosomal storage diseases (LSDs) are a group of over 50 inherited metabolic disorders caused by deficiency of specific lysosomal hydrolases, resulting in progressive accumulation of undegraded substrates within lysosomes. Gaucher disease is caused by glucocerebrosidase deficiency, presenting with hepatosplenomegaly, bone pain (Erlenmeyer flask deformity of the distal femur), anemia, thrombocytopenia, and Gaucher cells; enzyme replacement therapy and substrate reduction therapy are effective. Niemann-Pick disease is caused by sphingomyelinase deficiency (types A and B) causing accumulation of sphingomyelin in macrophages; Type A presents with severe neurovisceral disease and a cherry-red spot on the macula. Fabry disease is an X-linked $\alpha$-galactosidase A deficiency causing accumulation of globotriaosylceramide, presenting with acroparesthesias, angiokeratomas, corneal verticillata, hypohidrosis, progressive renal failure, hypertrophic cardiomyopathy, and stroke; enzyme replacement therapy and oral chaperone therapy slow disease progression. Tay-Sachs disease is caused by hexosaminidase A deficiency, presenting in infancy with progressive neurodegeneration, hypotonia, exaggerated startle response, cherry-red spot, and blindness. Metachromatic leukodystrophy is caused by arylsulfatase A deficiency, causing accumulation of sulfatides with progressive demyelination.

\subsection{Maple Syrup Urine Disease}
\label{sec:Maple Syrup Urine Disease}
Maple syrup urine disease (MSUD) is an autosomal recessive disorder caused by deficiency of the branched-chain $\alpha$-ketoacid dehydrogenase (BCKD) enzyme complex, leading to accumulation of branched-chain amino acids (BCAAs)---leucine, isoleucine, valine---and their corresponding ketoacids. Incidence is approximately 1 in 185,000 live births worldwide, but is higher in certain populations (Mennonite communities, 1 in 380). The condition results from deficiency of the BCKD enzyme complex. Clinical features include the classic (severe) form presenting in the first week of life with poor feeding, lethargy, vomiting, hypertonicity, opisthotonos, seizures, and a characteristic sweet, maple syrup odor of cerumen and urine; progressive neurologic deterioration leads to coma and respiratory failure if untreated. Diagnosis is by newborn screening (elevated leucine/isoleucine) and plasma amino acid analysis showing markedly elevated leucine, isoleucine, valine, and pathognomonic alloisoleucine. Management of acute decompensation includes aggressive fluid resuscitation, high-dose dextrose and insulin to promote anabolism, and dialysis for rapid leucine clearance; chronic management involves dietary restriction of BCAAs, thiamine supplementation (in responsive variants), and vigilant monitoring during intercurrent illness; liver transplantation provides a functional cure. Prognosis is favorable with early treatment.

\subsection{Mucopolysaccharidoses}
\label{sec:Mucopolysaccharidoses}
Mucopolysaccharidoses (MPS) are a group of lysosomal storage disorders caused by deficiency of enzymes required for the degradation of glycosaminoglycans (GAGs), leading to progressive accumulation in tissues. Hurler syndrome (MPS I-H, severe) is caused by $\alpha$-L-iduronidase deficiency, presenting with coarse facies, corneal clouding, hepatosplenomegaly, dysostosis multiplex, cardiac valvular disease, and severe neurodegeneration. Hunter syndrome (MPS II, X-linked) is caused by iduronate-2-sulfatase deficiency, similar to MPS I but without corneal clouding. Sanfilippo syndrome (MPS III) is primarily neurodegenerative with severe behavioral disturbance, sleep disturbance, progressive intellectual decline, and seizures. Diagnosis is by urinary GAG screening with fractionation, enzyme assay, and genetic testing. Treatment includes hematopoietic stem cell transplantation for MPS I-H before age 2.5 years and enzyme replacement therapy for MPS I, II, and IVA. Prognosis is guarded; enzyme replacement improves somatic features but does not cross the blood-brain barrier.

\subsection{Organic Acidemias}
\label{sec:Organic Acidemias}
Organic acidemias are a group of autosomal recessive disorders of amino acid, fatty acid, and carbohydrate catabolism causing accumulation of organic acids in blood and urine. Methylmalonic acidemia (MMA) is caused by deficiency of methylmalonyl-CoA mutase or cobalamin metabolism, presenting in infancy with recurrent episodes of metabolic acidosis, hyperammonemia, hypoglycemia, lethargy, vomiting, hypotonia, and failure to thrive; long-term complications include chronic kidney disease, pancreatitis, and neurologic deficits. Propionic acidemia (PA) is caused by propionyl-CoA carboxylase deficiency, presenting similarly with severe metabolic acidosis, ketosis, hyperammonemia, pancytopenia, and basal ganglia infarction. Isovaleric acidemia (IVA) is caused by isovaleryl-CoA dehydrogenase deficiency, producing a characteristic ``sweaty feet'' odor (isovaleric acid) and presenting with vomiting, metabolic acidosis, lethargy, and pancytopenia. Diagnosis of all organic acidemias is by urine organic acid analysis, plasma acylcarnitine profile, and genetic testing. Management includes protein restriction, carnitine supplementation, and specific therapies (metronidazole for MMA, glycine and carnitine for IVA). Prognosis is variable.

\subsection{Phenylketonuria}
\label{sec:Phenylketonuria}
Phenylketonuria (PKU) is an autosomal recessive disorder of phenylalanine metabolism caused by deficiency of phenylalanine hydroxylase (PAH), the enzyme that converts phenylalanine to tyrosine. The incidence is approximately 1 in 10,000--15,000 live births. The condition results from accumulation of phenylalanine and its metabolites in blood and tissues, causing impaired brain development and function. Clinical features include progressive developmental delay, intellectual disability, microcephaly, seizures, behavioral problems, a characteristic musty odor of urine and skin, and hypopigmentation (fair skin, light hair, blue eyes) due to impaired melanin synthesis. Newborns are asymptomatic at birth. Diagnosis is by newborn screening (elevated phenylalanine greater than 120 \textmu mol/L) confirmed by elevated blood phenylalanine and normal or low tyrosine. Management involves lifelong dietary restriction of phenylalanine (elimination of high-protein foods and use of phenylalanine-free medical formulas); sapropterin (BH4 supplementation) is effective in a subset of patients; pegvaliase is an enzyme substitution therapy for adults with poorly controlled PKU. Prognosis is excellent with early and consistent dietary treatment.

\subsection{Porphyrias}
\label{sec:Porphyrias}
The porphyrias are a group of inherited or acquired disorders of heme biosynthesis caused by deficiency of specific enzymes in the heme biosynthetic pathway, leading to accumulation of porphyrins or their precursors. They are classified as hepatic or erythropoietic. Acute intermittent porphyria (AIP) is caused by porphobilinogen deaminase deficiency (autosomal dominant with low penetrance), the most common acute hepatic porphyria; attacks are precipitated by drugs, fasting, infections, and hormonal changes. Clinical features include severe abdominal pain (90\%), neuropsychiatric symptoms, autonomic dysfunction, and peripheral motor neuropathy; cutaneous photosensitivity is absent. Diagnosis is by elevated urinary porphobilinogen (PBG) and aminolevulinic acid (ALA) during attacks. Management includes removing precipitating factors, high-carbohydrate loading, intravenous hemin, and givosiran for prevention. Porphyria cutanea tarda (PCT) is caused by uroporphyrinogen decarboxylase deficiency, the most common porphyria, associated with alcohol, hepatitis C, estrogen therapy, and hemochromatosis; clinical features include blistering skin lesions on sun-exposed areas; treatment includes phlebotomy or low-dose hydroxychloroquine. Erythropoietic protoporphyria (EPP) is caused by ferrochelatase deficiency, causing photosensitivity in childhood. Prognosis is variable.

\subsection{Urea Cycle Disorders}
\label{sec:Urea Cycle Disorders}
Urea cycle disorders (UCDs) are a group of inborn errors of metabolism caused by deficiency of one of the six enzymes or two transporters involved in hepatic conversion of ammonia to urea. X-linked ornithine transcarbamylase (OTC) deficiency is the most common (incidence 1 in 14,000). The condition results from impaired hepatic conversion of ammonia to urea, leading to hyperammonemia. Clinical features include severe enzyme deficiencies presenting in neonates after 24--72 hours with poor feeding, vomiting, lethargy, tachypnea with respiratory alkalosis, hypothermia, seizures, and rapidly progressive coma; late-onset forms present with recurrent hyperammonemia triggered by protein loads or catabolic stress, manifesting as developmental delay, ataxia, behavioral changes, and headaches. Diagnosis is by plasma ammonia (elevated, often greater than 200 \textmu mol/L), plasma amino acids, urinary orotic acid, and genetic testing. Acute management includes discontinuing protein intake, high-calorie dextrose and lipids to promote anabolism, nitrogen scavengers (IV sodium benzoate and sodium phenylacetate), arginine supplementation, and hemodialysis for severe hyperammonemia. Chronic management includes protein restriction, essential amino acid supplementation, oral nitrogen scavengers, and consideration of orthotopic liver transplantation. Prognosis depends on severity and prompt treatment.

%-------------------------------------------------------------
\section{Mineral Deficiency and Excess Disorders}
%-------------------------------------------------------------

%-------------------------------------------------------------

%-------------------------------------------------------------

%-------------------------------------------------------------

%-------------------------------------------------------------

%-------------------------------------------------------------

\label{sec:Mineral Deficiency and Excess Disorders}
%-------------------------------------------------------------%-------------------------------------------------------------

\subsection{Calcium and Phosphorus Disorders}
\label{sec:Calcium and Phosphorus Disorders}
Calcium and phosphorus disorders encompass a group of mineral disorders resulting from dysregulation of calcium and phosphorus homeostasis, which is tightly regulated by parathyroid hormone (PTH), vitamin D, and calcitonin. Hypocalcemia (total calcium less than 8.5 mg/dL or ionized calcium less than 4.2 mg/dL) results from vitamin D deficiency, hypoparathyroidism, chronic kidney disease, hypomagnesemia, and medications; clinical features include neuromuscular irritability (perioral paresthesias, Chvostek sign, Trousseau sign, carpopedal spasm, tetany), prolonged QT interval, seizures, and laryngospasm. Hypercalcemia (total calcium greater than 10.5 mg/dL) is most commonly due to primary hyperparathyroidism or malignancy; clinical features include ``bones, stones, groans, and psychiatric overtones''---bone pain, nephrolithiasis, polyuria, polydipsia, constipation, nausea, confusion, and shortened QT interval; severe hypercalcemia (greater than 13 mg/dL) is a medical emergency requiring IV fluids, calcitonin, bisphosphonates, and treatment of the underlying cause. Hypophosphatemia (phosphate less than 2.5 mg/dL) causes muscle weakness, rhabdomyolysis, and impaired bone mineralization; hyperphosphatemia (phosphate greater than 4.5 mg/dL) is most common in chronic kidney disease, causing secondary hyperparathyroidism, vascular calcification, and renal osteodystrophy.

\subsection{Copper Deficiency and Excess}
\label{sec:Copper Disorders}
Copper deficiency is a rare mineral deficiency disorder that occurs after gastric bypass surgery, in malabsorption, prolonged parenteral nutrition, excessive zinc intake (zinc induces metallothionein which blocks intestinal copper absorption), and in Menkes disease. The condition results from deficiency of copper, an essential trace element serving as a cofactor for cytochrome c oxidase, superoxide dismutase, lysyl oxidase, dopamine $\beta$-hydroxylase, and ceruloplasmin. Clinical features include microcytic hypochromic anemia (refractory to iron therapy), neutropenia, myeloneuropathy (spastic gait, sensory ataxia, peripheral neuropathy mimicking B12 deficiency), and pancytopenia. Diagnosis is by low serum copper and low ceruloplasmin. Management involves oral copper 2--5 mg/kg/day as copper sulfate. Menkes disease (kinky hair disease) is an X-linked recessive disorder of ATP7A causing impaired intestinal copper absorption, presenting in infancy with progressive neurodegeneration, seizures, hypotonia, failure to thrive, characteristic brittle, depigmented, kinky hair (pili torti), and tortuous cerebral arteries; subcutaneous copper-histidine may improve outcomes if started early, but prognosis is poor. Wilson disease (copper excess) is covered in Chapter~8.

\subsection{Fluoride Excess --- Fluorosis}
\label{sec:Fluorosis}
Fluorosis is a mineral excess disorder caused by chronic excessive intake of fluoride, typically from drinking water with high fluoride content (greater than 1.5 ppm), but also from dental products, fluoride supplements, and industrial exposure. It is endemic in regions of India, China, Africa, and the Middle East. The condition results from fluoride incorporation into hydroxyapatite crystals of bones and teeth. Clinical features include dental fluorosis (chalky white or yellow-brown mottling, pitting, and staining of teeth, occurring when excess fluoride is ingested during tooth development) and skeletal fluorosis (osteosclerosis, bone and joint pain, ligamentous calcification, periosteal bone formation, and spinal rigidity mimicking ankylosing spondylitis; advanced stages lead to crippling deformities, kyphosis, and neurologic complications from spinal canal stenosis). Diagnosis is based on history of high fluoride exposure, radiographic findings (osteosclerosis, periosteal bone apposition, interosseous membrane calcification), and elevated urinary fluoride. Management focuses on prevention: defluoridation of drinking water, use of alternative water sources, and avoidance of fluoride-containing products in endemic areas. Prognosis is good with prevention; no specific treatment exists for established disease.

\subsection{Iodine Deficiency Disorders}
\label{sec:Iodine Deficiency Disorders}
Iodine deficiency disorders are a group of mineral deficiency disorders affecting nearly 2 billion people, particularly in mountainous regions (Himalayas, Andes, Alps), flood-prone river valleys, and areas with iodine-depleted soils. The condition results from deficiency of iodine, an essential component of thyroid hormones (thyroxine, T4; triiodothyronine, T3) required for normal growth, development, and metabolism. Clinical features include goiter (compensatory thyroid enlargement), cretinism (neurologic cretinism with intellectual disability, deaf-mutism, spastic diplegia; myxedematous cretinism with hypothyroidism, growth retardation, coarse facial features), hypothyroidism in older children and adults (fatigue, cold intolerance, weight gain, constipation, dry skin, bradycardia), and pregnancy complications including miscarriage, stillbirth, congenital anomalies, and impaired neurocognitive development. Diagnosis is by low urinary iodine concentration (less than 100 \textmu g/L), elevated TSH, low T4, and goiter on palpation or ultrasound. Management involves universal salt iodization programs for prevention; iodine supplementation (oral potassium iodide 100--200 mg daily) for mild deficiency and levothyroxine for hypothyroidism. Prognosis is excellent with correction of deficiency.

\subsection{Iron Deficiency}
\label{sec:Iron Deficiency}
Iron deficiency is the most common nutritional deficiency worldwide and the leading cause of anemia, affecting approximately 25\% of the global population. The condition results from inadequate dietary iron intake, poor bioavailability (non-heme iron from plant sources), increased demands (pregnancy, growth, menstruation), and chronic blood loss (gastrointestinal bleeding, menorrhagia, hookworm); iron is essential for hemoglobin synthesis, myoglobin, and numerous enzymes involved in electron transport and DNA synthesis. Clinical features manifest as microcytic hypochromic anemia including fatigue, pallor, weakness, dyspnea on exertion, tachycardia, headache, and epithelial changes including koilonychia (spoon nails), glossitis, angular stomatitis, and pagophagia (pica for ice). Diagnosis is by low serum ferritin (gold standard), low serum iron, low transferrin saturation, elevated total iron-binding capacity (TIBC), elevated soluble transferrin receptor, and peripheral smear showing microcytic, hypochromic RBCs with pencil-shaped cells and target cells. Management involves oral iron therapy (ferrous sulfate 325 mg three times daily, equivalent to 195 mg elemental iron per day) as first-line, with response monitored by reticulocytosis at 3--7 days and hemoglobin rise of 1 g/dL per 2--3 weeks; parenteral iron is used for intolerance, severe deficiency, malabsorption, or chronic kidney disease; treatment of the underlying cause is mandatory. Prognosis is excellent with treatment.

\subsection{Selenium Deficiency}
\label{sec:Selenium Deficiency}
Selenium deficiency is a mineral deficiency disorder occurring in regions with low soil selenium content (parts of China, Europe, New Zealand). The condition results from deficiency of selenium, an essential trace element functioning as a component of selenoproteins, including glutathione peroxidases (antioxidant defense), iodothyronine deiodinases (thyroid hormone metabolism), and selenoprotein P (selenium transport). Clinical features include Keshan disease (an endemic dilated cardiomyopathy affecting children and women of childbearing age, presenting with acute or chronic heart failure, cardiomegaly, arrhythmias) and Kashin-Beck disease (an osteoarthropathy presenting with joint pain, stiffness, and deformities of the fingers, knees, and ankles); selenium deficiency also impairs immune function, exacerbates viral virulence, and is associated with infertility and myopathy. Diagnosis is by low serum selenium (less than 50--70 \textmu g/L) and decreased glutathione peroxidase activity. Management involves sodium selenite or selenomethionine 100--200 \textmu g daily. Prognosis is variable.

\subsection{Zinc Deficiency}
\label{sec:Zinc Deficiency}
Zinc deficiency is a mineral deficiency disorder common in low-resource settings due to low dietary intake of animal-source foods and high phytate content (inhibits zinc absorption). At-risk groups include pregnant and lactating women, children (especially with diarrhea), the elderly, vegetarians, and those with malabsorption (acrodermatitis enteropathica---autosomal recessive zinc transporter defect). The condition results from deficiency of zinc, an essential trace element that serves as a cofactor for over 300 enzymes involved in growth, immune function, protein synthesis, wound healing, and DNA synthesis. Clinical features include growth retardation, delayed sexual maturation, hypogonadism, impaired immune function (recurrent infections, especially diarrheal diseases), alopecia, dermatitis (periorificial and acral, vesiculobullous or eczematous---characteristic of acrodermatitis enteropathica), poor wound healing, night blindness, anorexia, dysgeusia, and neuropsychiatric changes. Diagnosis is by low plasma zinc (less than 60--70 \textmu g/dL). Management involves zinc sulfate or acetate 1--3 mg/kg/day (adults: 25--50 mg elemental zinc daily); for acrodermatitis enteropathica, lifelong supplementation is required. Prognosis is excellent with supplementation.

%-------------------------------------------------------------
\section{Obesity and Overweight}
%-------------------------------------------------------------

%-------------------------------------------------------------

%-------------------------------------------------------------

%-------------------------------------------------------------

%-------------------------------------------------------------

%-------------------------------------------------------------

\label{sec:Obesity and Overweight}
%-------------------------------------------------------------%-------------------------------------------------------------

\subsection{Metabolic Syndrome}
\label{sec:Metabolic Syndrome}
Metabolic syndrome is a cluster of interconnected metabolic abnormalities that increase the risk of cardiovascular disease, type 2 diabetes, and all-cause mortality. The NCEP ATP III diagnostic criteria require at least three of the following: central obesity (waist circumference greater than 102 cm in men, 88 cm in women), elevated triglycerides (150 mg/dL or greater), reduced HDL cholesterol (less than 40 mg/dL in men, less than 50 mg/dL in women), elevated blood pressure (130/85 mm Hg or greater), and elevated fasting glucose (100 mg/dL or greater). Central obesity and insulin resistance are considered core pathophysiologic features, resulting from adipose tissue dysfunction and ectopic fat deposition driving insulin resistance, endothelial dysfunction, and a prothrombotic/proinflammatory state. Clinical features include the cluster of metabolic abnormalities without specific symptoms. Diagnosis requires fulfillment of the diagnostic criteria. Management focuses on lifestyle modification and pharmacologic treatment of individual risk factors. Prognosis is variable; metabolic syndrome confers a 2-fold increased risk of cardiovascular events and a 5-fold increased risk of developing type 2 diabetes.

\subsection{Obesity}
\label{sec:Obesity}
Obesity is a chronic, relapsing, multifactorial disease characterized by excessive accumulation of adipose tissue with adverse health consequences. It is defined by a body mass index (BMI) of 30 kg/m\textsuperscript{2} or greater in adults, classified into class I (30--34.9), class II (35--39.9), and class III (40 or higher). The global prevalence has nearly tripled since 1975; over 650 million adults are obese. The condition results from a chronic energy imbalance---caloric intake exceeding energy expenditure---driven by complex interactions of genetic, environmental, behavioral, psychosocial, endocrine, and iatrogenic factors; adipose tissue is an active endocrine organ secreting adipokines (leptin, adiponectin, resistin, TNF$\alpha$, IL-6) contributing to a chronic low-grade inflammatory state. Clinical consequences affect virtually every organ system: type 2 diabetes mellitus, hypertension, dyslipidemia, coronary artery disease, stroke, obstructive sleep apnea, non-alcoholic fatty liver disease, cholelithiasis, osteoarthritis, infertility, depression, and increased risk of multiple cancers. Diagnosis includes BMI measurement, waist circumference, and screening for comorbidities. Management is stepwise: lifestyle intervention (dietary modification, physical activity, behavioral therapy), pharmacotherapy (orlistat, GLP-1 receptor agonists such as semaglutide and liraglutide, naltrexone-bupropion, phentermine-topiramate), and bariatric surgery for class III obesity or class II with comorbidities. Prognosis is variable; sustained weight loss improves or resolves many comorbidities.

\subsection{Sarcopenic Obesity}
\label{sec:Sarcopenic Obesity}
Sarcopenic obesity is a distinct clinical entity characterized by the combination of excess adiposity and low muscle mass and function in the same individual. It is increasingly recognized with aging, but also occurs in the setting of malignancy, chronic disease, and catabolic states. The condition results from a complex interplay of age-related hormonal decline, chronic inflammation, insulin resistance, physical inactivity, and inadequate protein intake. Clinical features include the combination of obesity with low muscle strength and functional decline. Diagnosis requires measurement of body composition (DXA, bioelectrical impedance analysis, CT, or MRI) along with assessment of muscle strength and physical performance. Management includes resistance exercise training, adequate protein intake, vitamin D supplementation, and intentional weight loss through caloric restriction with emphasis on preserving lean mass. Prognosis is guarded; sarcopenic obesity carries additive risks for falls, fractures, physical disability, cardiovascular disease, and mortality.

%-------------------------------------------------------------
\section{Protein-Energy Malnutrition}
%-------------------------------------------------------------

%-------------------------------------------------------------

%-------------------------------------------------------------

%-------------------------------------------------------------

%-------------------------------------------------------------

%-------------------------------------------------------------

\label{sec:Protein-Energy Malnutrition}
%-------------------------------------------------------------%-------------------------------------------------------------

Protein-energy malnutrition (PEM) encompasses a spectrum of disorders arising from inadequate intake of calories, protein, or both. It remains a critical public health problem in low- and middle-income countries, contributing to approximately 45\% of deaths in children under five years globally. PEM increases susceptibility to infections, impairs growth and development, and has lifelong consequences for physical and cognitive function. Climate-related food insecurity, conflict, and poverty are major drivers of its persistence.

\subsection{Kwashiorkor}
\label{sec:Kwashiorkor}
Kwashiorkor is a form of severe protein-energy malnutrition resulting from inadequate protein intake despite relatively preserved caloric intake. It typically occurs in children aged 1--3 years after weaning, particularly when the diet consists predominantly of starchy, low-protein foods. The condition results from hypoalbuminemia due to insufficient dietary protein, leading to decreased plasma oncotic pressure and peripheral edema. Clinical features include bilateral pitting edema (beginning in the feet and ascending), hepatomegaly (fatty liver), thinning and depigmentation of hair (flag sign), a characteristic ``flaky paint'' dermatitis, cheilosis, and stomatitis; affected children appear miserable and apathetic. The presence of edema distinguishes kwashiorkor from marasmus. Diagnosis is clinical, supplemented by low serum albumin (less than 25 g/L). Management follows the WHO protocol: treatment of hypoglycemia and hypothermia, empiric antibiotics, correction of electrolyte imbalances, and staged refeeding; transition from F-75 to F-100 formula is guided by clinical response. Prognosis is good with treatment; edema usually resolves within 1--2 weeks, but long-term cognitive deficits may persist.

\subsection{Marasmic-Kwashiorkor}
\label{sec:Marasmic-Kwashiorkor}
Marasmic-kwashiorkor is a mixed form of severe acute malnutrition combining features of both marasmus (wasting) and kwashiorkor (edema). It occurs when there is deficiency of both total calories and protein. Clinical features include severe wasting accompanied by peripheral edema. Anthropometry shows severe deficits in weight-for-age and weight-for-height, with concurrent pitting edema and low serum albumin. Diagnosis is clinical with anthropometric confirmation. Management follows the same WHO guidelines with emphasis on careful monitoring for fluid overload, refeeding syndrome, and infections. Prognosis is worse than for either pure form alone.

\subsection{Marasmus}
\label{sec:Marasmus}
Marasmus is a form of severe protein-energy malnutrition characterized by inadequate intake of both protein and calories. It is most common in infants and young children in low-resource settings, particularly during the first year of life, and is associated with early weaning, repeated infections, and poverty. The condition results from catabolism of muscle and subcutaneous fat to meet energy needs, leading to profound wasting with relative preservation of serum albumin. Clinical features include severe wasting of muscle and subcutaneous tissue, a ``monkey-like'' facies with prominent eyes, wrinkled skin, irritability, and growth retardation; subcutaneous fat is absent and there is no peripheral edema. Diagnosis is based on anthropometric criteria: weight-for-age less than 60\% of median, weight-for-height less than 70\% of median, or mid-upper arm circumference (MUAC) less than 115 mm in children aged 6--59 months. Management follows the WHO protocol: stabilization with correction of dehydration and electrolyte imbalances (using ReSoMal), empiric antibiotics, micronutrient supplementation, and staged refeeding with F-75 formula followed by F-100 formula; slow, cautious refeeding is required to avoid refeeding syndrome. Prognosis is good with timely treatment, but mortality remains high (20--30\%) in severe cases without adequate care.

%-------------------------------------------------------------
\section{Vitamin Deficiency Disorders}
%-------------------------------------------------------------

%-------------------------------------------------------------

%-------------------------------------------------------------

%-------------------------------------------------------------

%-------------------------------------------------------------

%-------------------------------------------------------------

\label{sec:Vitamin Deficiency Disorders}
%-------------------------------------------------------------%-------------------------------------------------------------

\subsection{Vitamin A Deficiency}
\label{sec:Vitamin A Deficiency}
Vitamin A deficiency (VAD) is a vitamin deficiency disorder that is the leading cause of preventable blindness in children worldwide and a major contributor to morbidity from infections. It is most prevalent in sub-Saharan Africa and South Asia, where diets are low in preformed vitamin A (retinol, retinyl esters) and provitamin A carotenoids. The condition results from inadequate dietary intake of vitamin A, which is essential for retinal photoreceptor function (11-\emph{cis}-retinal in rhodopsin), epithelial integrity, and immune competence. Clinical features include night blindness (nyctalopia) as the earliest symptom, followed by xerophthalmia including conjunctival xerosis, Bitot spots, corneal xerosis, and ultimately keratomalacia causing irreversible blindness; VAD also increases the severity of measles, diarrheal diseases, and respiratory infections. Diagnosis is clinical (ocular signs) and confirmed by serum retinol less than 0.70 \textmu mol/L. Management involves high-dose vitamin A supplementation (200,000 IU orally for children over 12 months, 100,000 IU for 6--12 months, 50,000 IU for infants under 6 months) given on diagnosis, repeated the next day, and at least 2 weeks later; prevention includes dietary diversification, food fortification, and periodic high-dose supplementation. Prognosis is excellent for early stages; corneal involvement may cause irreversible blindness.

\subsection{Vitamin B1 (Thiamine) Deficiency --- Beriberi and Wernicke-Korsakoff Syndrome}
\label{sec:Vitamin B1 Deficiency}
Thiamine (vitamin B1) deficiency is a vitamin deficiency disorder affecting populations consuming polished white rice as a dietary staple, individuals with chronic alcoholism, hyperemesis gravidarum, after bariatric surgery, and in prolonged parenteral nutrition without thiamine supplementation. The condition results from thiamine deficiency; thiamine functions as a cofactor for several enzymes involved in carbohydrate metabolism, including transketolase, pyruvate dehydrogenase, and $\alpha$-ketoglutarate dehydrogenase. Clinical features include two major syndromes: dry beriberi (symmetric peripheral neuropathy with stocking-glove sensory loss, muscle weakness, absent deep tendon reflexes, paresthesias), wet beriberi (high-output heart failure with dilated cardiomyopathy, peripheral vasodilation, edema, tachycardia, warm extremities, bounding pulses, lactic acidosis), and Wernicke encephalopathy (acute neuropsychiatric emergency with confusion, ataxia, and ophthalmoplegia) which may progress to Korsakoff syndrome (irreversible amnestic disorder with anterograde and retrograde amnesia and confabulation). Diagnosis is clinical; whole blood thiamine or erythrocyte transketolase activity may be measured but treatment should not be delayed. Management involves thiamine 100--500 mg IV/IM daily for 5--7 days followed by oral maintenance; Wernicke encephalopathy requires immediate high-dose parenteral thiamine before glucose administration. Prognosis is favorable for beriberi with treatment; Wernicke encephalopathy is reversible with prompt treatment, but Korsakoff syndrome is irreversible.

\subsection{Vitamin B12 (Cobalamin) Deficiency}
\label{sec:Vitamin B12 Deficiency}
Vitamin B12 (cobalamin) deficiency is a vitamin deficiency disorder common in the elderly due to atrophic gastritis and declining intrinsic factor production (food-cobalamin malabsorption). Pernicious anemia is an autoimmune disorder with antibodies against gastric parietal cells and/or intrinsic factor, leading to profound B12 deficiency. Other causes include gastric or ileal resection, Crohn disease, bacterial overgrowth, fish tapeworm (Diphyllobothrium latum) infection, strict vegetarian/vegan diet, and metformin use. The condition results from deficiency of vitamin B12, which is essential for DNA synthesis, myelin formation, and homocysteine metabolism. Clinical features include megaloblastic anemia and the hallmark distinguishing feature from folate deficiency: neurologic involvement (subacute combined degeneration of the spinal cord manifesting as symmetric paresthesias, loss of vibration and proprioception, ataxia, spastic paraparesis, and a positive Romberg sign), peripheral neuropathy, visual disturbances (optic atrophy), memory loss, and psychiatric symptoms. Diagnosis is by low serum B12 (less than 200 pg/mL) with elevated serum methylmalonic acid (MMA) and homocysteine; elevated MMA is the most sensitive indicator. Management involves oral B12 1000--2000 \textmu g daily for deficiency without neurologic involvement; for pernicious anemia, neurologic symptoms, or severe deficiency, parenteral B12 (hydroxocobalamin or cyanocobalamin) 1000 \textmu g IM daily for 1 week, weekly for 4 weeks, then monthly for life. Prognosis is excellent with treatment; hematologic response is rapid (reticulocytosis within 3--5 days, hemoglobin normalization in 6--8 weeks).

\subsection{Vitamin B2 (Riboflavin) Deficiency --- Ariboflavinosis}
\label{sec:Vitamin B2 Deficiency}
Riboflavin (vitamin B2) deficiency (ariboflavinosis) is a vitamin deficiency disorder that is rare in developed countries but occurs in alcoholism, malabsorption syndromes, and populations with limited dairy and meat intake. The condition results from deficiency of riboflavin, a precursor for flavin mononucleotide (FMN) and flavin adenine dinucleotide (FAD), cofactors for numerous redox enzymes. Clinical features include angular stomatitis (cheilosis), glossitis (smooth, magenta-colored tongue), seborrheic dermatitis (nasolabial folds, scrotum, vulva), photophobia, corneal vascularization, and normocytic anemia. Diagnosis is clinical and confirmed by low erythrocyte glutathione reductase activity (EGRAC). Management involves riboflavin 5--10 mg orally daily. Prognosis is excellent; mucocutaneous lesions resolve promptly with treatment.

\subsection{Vitamin B3 (Niacin) Deficiency --- Pellagra}
\label{sec:Vitamin B3 Deficiency}
Niacin (vitamin B3) deficiency (pellagra) is a vitamin deficiency disorder classically associated with diets based on maize (corn), which contains niacin in a bound, non-bioavailable form unless treated with alkali (nixtamalization); it also occurs in alcoholism, carcinoid syndrome, and with isoniazid therapy. The condition results from deficiency of niacin, a precursor for nicotinamide adenine dinucleotide (NAD) and NADP, essential for redox reactions and cellular energy metabolism. Clinical features include the classic triad of dermatitis, dementia, and diarrhea (the ``three D's''); dermatitis is photosensitive, symmetric, with well-demarcated, erythematous, hyperpigmented, scaling plaques in sun-exposed areas (Casal necklace); gastrointestinal involvement includes diarrhea, glossitis, and stomatitis; neuropsychiatric manifestations include depression, apathy, confusion, memory loss, and psychosis. Diagnosis is clinical. Management involves nicotinamide (preferred over niacin to avoid flushing) 100--500 mg orally daily. Prognosis is excellent with treatment; dermatitis and gastrointestinal symptoms improve within days; neurologic symptoms resolve more slowly.

\subsection{Vitamin B6 (Pyridoxine) Deficiency}
\label{sec:Vitamin B6 Deficiency}
Pyridoxine (vitamin B6) deficiency is a vitamin deficiency disorder occurring in alcoholism, chronic renal failure, autoimmune disorders (rheumatoid arthritis), and with medications such as isoniazid, cycloserine, hydralazine, and oral contraceptives. The condition results from deficiency of pyridoxine, which functions as a cofactor (pyridoxal phosphate, PLP) for over 100 enzymes involved in amino acid metabolism, neurotransmitter synthesis, and heme biosynthesis. Clinical features include microcytic hypochromic anemia (sideroblastic anemia with ringed sideroblasts in bone marrow), seborrheic dermatitis, glossitis, cheilosis, peripheral neuropathy, depression, confusion, and in infants, irritability and seizures; homocysteine elevation due to B6 deficiency is a cardiovascular risk factor. Diagnosis is by low plasma PLP levels. Management involves pyridoxine 50--100 mg orally daily; isoniazid-induced deficiency requires 50--100 mg of pyridoxine daily as prophylaxis. Prognosis is excellent with treatment.

\subsection{Vitamin B7 (Biotin) Deficiency}
\label{sec:Vitamin B7 Deficiency}
Biotin (vitamin B7) deficiency is a rare vitamin deficiency disorder because biotin is widely distributed in foods and synthesized by gut microbiota. Causes include high consumption of raw egg whites (avidin binds biotin, preventing absorption), parenteral nutrition without biotin, prolonged antibiotic use, and biotinidase deficiency (an inborn error). The condition results from deficiency of biotin, a cofactor for carboxylase enzymes involved in gluconeogenesis, fatty acid synthesis, and amino acid catabolism. Clinical features include periorificial dermatitis, alopecia, brittle nails, glossitis, conjunctivitis, neurologic symptoms (depression, lethargy, hallucinations, paresthesias), and metabolic acidosis; in biotinidase deficiency, presentation is in infancy with seizures, hypotonia, ataxia, hearing loss, and developmental delay. Diagnosis is by low serum biotin or elevated urinary 3-hydroxyisovaleric acid. Management involves biotin 5--10 mg orally daily (or 10--40 mg daily in biotinidase deficiency). Prognosis is excellent with treatment.

\subsection{Vitamin B9 (Folate) Deficiency}
\label{sec:Vitamin B9 Deficiency}
Folate (vitamin B9) deficiency is one of the most common vitamin deficiencies worldwide. Causes include inadequate dietary intake (especially in the elderly, alcoholics, and those with poor vegetable intake), malabsorption (celiac disease, tropical sprue, inflammatory bowel disease), increased requirements (pregnancy, lactation, hemolytic anemias, malignancy), and medications (methotrexate, trimethoprim, sulfasalazine, phenytoin). The condition results from deficiency of folate, a water-soluble B vitamin essential for one-carbon transfer reactions, DNA synthesis, purine and pyrimidine synthesis, and homocysteine remethylation. Clinical features overlap with vitamin B12 deficiency and include megaloblastic anemia (macrocytic RBCs with oval macrocytes, hypersegmented neutrophils, pancytopenia), glossitis (smooth, beefy red tongue), diarrhea, and infertility; unlike B12 deficiency, folate deficiency does NOT cause neurologic manifestations. Diagnosis is confirmed by serum folate less than 2 ng/mL and elevated homocysteine with normal methylmalonic acid (key to distinguish from B12 deficiency). Management involves folic acid 1--5 mg orally daily; periconceptional folic acid supplementation (400 \textmu g/day) prevents neural tube defects. Prognosis is excellent with treatment; before treating macrocytic anemia, vitamin B12 deficiency must be ruled out or treated concurrently to prevent precipitating subacute combined degeneration.

\subsection{Vitamin C Deficiency --- Scurvy}
\label{sec:Vitamin C Deficiency}
Ascorbic acid (vitamin C) deficiency (scurvy) is a vitamin deficiency disorder that is rare in developed countries but occurs in individuals with severely restricted diets (alcoholics, elderly, psychiatric patients, food faddism), malabsorption, and in infants fed only cow's milk without supplementation. The condition results from deficiency of ascorbic acid, which is essential for collagen synthesis (hydroxylation of proline and lysine in procollagen), carnitine synthesis, neurotransmitter synthesis, and as an antioxidant. Clinical features reflect impaired collagen synthesis: perifollicular hemorrhages, petechiae, ecchymoses, gingival hyperplasia with bleeding, impaired wound healing, hyperkeratosis, coiled hairs (corkscrew hairs), and leg edema; musculoskeletal manifestations include arthralgias, myalgias, and subperiosteal hemorrhages (particularly painful in infants presenting with pseudoparalysis); systemic symptoms include fatigue, irritability, depression, and anemia. Diagnosis is clinical; serum ascorbic acid less than 0.2 mg/dL confirms deficiency. Management involves vitamin C 100--200 mg orally three times daily for 1 week, then 100 mg daily; infantile scurvy responds to 50 mg daily. Prognosis is excellent; clinical improvement occurs within days.

\subsection{Vitamin D Deficiency --- Rickets and Osteomalacia}
\label{sec:Vitamin D Deficiency}
Vitamin D deficiency is a vitamin deficiency disorder resulting from inadequate sunlight exposure, low dietary intake (or lack of fortified foods), malabsorption (celiac disease, cystic fibrosis, short bowel syndrome), renal disease (impaired 1$\alpha$-hydroxylation), obesity (sequestration in adipose tissue), and medications (antiepileptics, rifampin, glucocorticoids). The condition results from deficiency of vitamin D, a fat-soluble secosteroid essential for intestinal calcium absorption, bone mineralization, and calcium homeostasis. Clinical features in children include rickets (defective mineralization of growing bone) with craniotabes, frontal bossing, delayed fontanelle closure, rachitic rosary, Harrison groove, bowing of the legs (genu varum), widened wrists and ankles, myopathy, and hypocalcemic tetany. In adults, osteomalacia presents with bone pain, proximal muscle weakness (waddling gait), and fragility fractures. Diagnosis is by serum 25-hydroxyvitamin D less than 20 ng/mL (50 nmol/L), elevated alkaline phosphatase, low-normal calcium, and low phosphate; radiographs of rickets show cupping, fraying, and widening of the metaphysis; in osteomalacia, Looser zones (pseudofractures) are seen. Management involves vitamin D 2000--5000 IU daily for 12 weeks or 50,000 IU weekly for 8 weeks, followed by maintenance 800--2000 IU daily; correction of calcium and phosphate is essential. Prognosis is excellent with treatment.

\subsection{Vitamin E Deficiency}
\label{sec:Vitamin E Deficiency}
Vitamin E ($\alpha$-tocopherol) deficiency is a rare vitamin deficiency disorder occurring primarily in fat malabsorption syndromes (cystic fibrosis, cholestatic liver disease, abetalipoproteinemia, short bowel syndrome) and in genetic disorders of $\alpha$-tocopherol transfer protein (ataxia with isolated vitamin E deficiency, AVED). The condition results from deficiency of vitamin E, a fat-soluble antioxidant that protects cell membranes from lipid peroxidation. Clinical features include progressive neurologic dysfunction: spinocerebellar degeneration with ataxia, areflexia, loss of vibration and proprioception, ophthalmoplegia, and pigmentary retinopathy; hemolytic anemia may occur in premature infants. Diagnosis is by low plasma $\alpha$-tocopherol corrected for serum lipids. Management involves $\alpha$-tocopherol 800--1200 IU daily orally (or higher doses in severe malabsorption). Prognosis is variable; neurologic deficits may stabilize but are not fully reversible.

\subsection{Vitamin K Deficiency}
\label{sec:Vitamin K Deficiency}
Vitamin K deficiency is a vitamin deficiency disorder leading to impaired hepatic synthesis of functional clotting factors. Causes include inadequate dietary intake (rare in healthy adults), fat malabsorption, liver disease, antibiotic therapy (depleting gut flora), and parenteral nutrition. The condition results from deficiency of vitamin K, a fat-soluble vitamin required for the $\gamma$-carboxylation of glutamic acid residues on coagulation factors II (prothrombin), VII, IX, and X, as well as proteins C and S and bone matrix proteins. Clinical features include vitamin K deficiency bleeding (VKDB) in neonates presenting with bruising, gastrointestinal bleeding, umbilical bleeding, and life-threatening intracranial hemorrhage; in adults, deficiency presents with easy bruising, epistaxis, hematuria, and prolonged prothrombin time (PT) with normal PTT. Diagnosis is by prolonged PT (factor VII has the shortest half-life) that corrects with normal plasma, low factor II, VII, IX, X activity, and normal platelet count and fibrinogen. Management involves vitamin K1 (phylloquinone) 1--10 mg orally or 1--2 mg subcutaneously/IM for adults; VKDB prophylaxis includes 1 mg IM at birth; active bleeding requires fresh frozen plasma (FFP) or prothrombin complex concentrates (PCC) with parenteral vitamin K. Prognosis is excellent with treatment.

